\chapter*{Abstract}
\addcontentsline{toc}{chapter}{Abstract}

Despite over USD 1.3 trillion in announced subsidies and 422 GW of project pipeline, only 7\% of 2023 announced green hydrogen capacity reached scheduled Final Investment Decision \citep{odenweller2025green}. This thesis provides the first cross-jurisdictional causal evaluation of why low-carbon hydrogen projects survive or fail, and which policy mechanisms move the needle.

Using a project-level dataset of 1,354 Blue and Green hydrogen projects \mbox{(S\&P Global, 2010--2024)} with 367 documented failures, this research makes four contributions. \emph{First}, it identifies a previously under-appreciated mechanism: pre-FID offtake commitments reduce project-failure probability by 11--13 percentage points across five independent identification strategies, with Oster (2019) selection-on-unobservables sensitivity $\delta_{\text{null}} = 20.23$ indicating exceptional robustness. \emph{Second}, it evaluates four carrot-policy types (US Section 45Q, EU Innovation Fund, UK Track-1, China 14th Five-Year Plan) via modern difference-in-differences estimators \citep{sun2021estimating,borusyak2024revisiting}, with Rambachan-Roth (2023) honest sensitivity bounds revealing heterogeneous causal-identification strength: only China's state-coordinated mandate survives the strictest robustness checks ($M^* = 1.5$), while US, EU, and UK estimates are sensitivity-bounded. \emph{Third}, it documents a structural break in the policy-effect coefficient around 2020 via three time-varying-parameter estimators (threshold, AR(1) state-space, random walk), consistent with post-announcement-boom speculative entry. \emph{Fourth}, it grounds these findings in a Dixit-Pindyck (1994) real-options framework that distinguishes $V/I$-boost mechanisms (output credits, capex grants) from $\sigma$-attack mechanisms (offtake mandates, cluster tenders).

Counterfactual scenario analysis quantifies external validity: an EU sector-optimal carrot mix would deliver an estimated +113 additional FIDs, +7.83 Mt/y additional CO\textsubscript{2} capture, and +1.76 Mt/y additional H\textsubscript{2} output over a three-year horizon (95\% bootstrap CI: +82 to +141 FIDs). These findings are externally corroborated by the 2024--2026 wave of high-profile cancellations: ArcelorMittal's withdrawal from Bremen and Eisenhüttenstadt despite \euro 1.3B in committed EU Innovation Fund subsidies, BP's exits from HyGreen Teesside and H2Teesside, and the September 2025 withdrawal of seven EU Hydrogen Bank auction winners representing 1.88 GW of combined capacity.

\bigskip
\noindent\textbf{Keywords:} hydrogen, causal inference, difference-in-differences, real options, mechanism design, policy evaluation, survival analysis, sensitivity analysis.

\noindent\textbf{JEL classification:} C21, C23, Q42, Q48, Q54, H23.

\clearpage
